\documentclass[11pt,a4paper]{article}

% ============================================================
% Conduite de projet : décodage quantique par réseau de neurones sur graphes
% Fichier autonome ; compilation recommandée avec pdfLaTeX.
% ============================================================
\usepackage[utf8]{inputenc}
\usepackage[T1]{fontenc}
\usepackage[french]{babel}
\usepackage{lmodern}
\usepackage{microtype}
\usepackage[a4paper,margin=2.35cm]{geometry}
\usepackage{amsmath,amssymb,mathtools}
\usepackage{booktabs,tabularx,longtable,array}
\usepackage{xcolor}
\usepackage{enumitem}
\usepackage{hyperref}
\usepackage{listings}
\usepackage{graphicx}
\usepackage{float}

\definecolor{navy}{HTML}{0C2340}
\definecolor{cyan}{HTML}{159E9C}
\definecolor{lightgray}{HTML}{F4F6F8}
\hypersetup{
  colorlinks=true,
  linkcolor=navy,
  citecolor=cyan,
  urlcolor=cyan,
  pdftitle={Conduite de projet -- Décodage quantique par GNN},
  pdfauthor={}
}
\setlist[itemize]{topsep=3pt,itemsep=2pt,leftmargin=1.5em}
\setlist[enumerate]{topsep=3pt,itemsep=3pt,leftmargin=1.8em}
\lstset{
  basicstyle=\ttfamily\small,
  breaklines=true,
  frame=single,
  backgroundcolor=\color{lightgray},
  keywordstyle=\color{navy},
  commentstyle=\color{cyan!70!black},
  showstringspaces=false
}
\newcommand{\pL}{p_{\mathrm L}}
\newcommand{\pphys}{p}
\newcommand{\GNN}{\mathrm{GNN}}
\newcommand{\PM}{\mathrm{PM}}
\newcommand{\xor}{\oplus}
\newcommand{\R}{\mathbb{R}}
\newcommand{\E}{\mathbb{E}}
\newcommand{\Prb}{\mathbb{P}}

\title{\textcolor{navy}{Conduite de projet}\\[0.4em]
\Large Décodage d'erreurs quantiques par réseau de neurones sur graphes}
\author{}
\date{Version de travail -- octobre 2026}

\begin{document}
\maketitle
\tableofcontents
\newpage

% ============================================================
\section{Résumé exécutif et objectif du projet}
% ============================================================

Ce projet étudie un problème concret de correction d'erreurs quantiques :
à partir des résultats de mesures de contrôle effectuées pendant une
mémoire quantique bruitée, déterminer si l'information logique a subi une
erreur. Le projet compare deux méthodes de décodage :
\begin{itemize}
  \item \textbf{PyMatching}, un décodeur classique par appariement de poids
  minimal (\emph{minimum-weight perfect matching}, MWPM), qui exploite un
  modèle de bruit fourni ;
  \item \textbf{un réseau de neurones sur graphes} (\emph{graph neural
  network}, GNN), entraîné sur des exemples simulés étiquetés.
\end{itemize}

\paragraph{Question scientifique principale.}
Pour des codes de surface de distances $d=3$ et $d=5$, sous un modèle de
bruit dépolarisant au niveau du circuit, quelle différence de taux d'erreur
logique observe-t-on entre un GNN entraîné sur des exemples simulés et
PyMatching, lorsque les deux méthodes sont évaluées sur les mêmes tirages
de test indépendants ? Quel est le coût de calcul associé à chaque méthode ?

\paragraph{Extension ciblée.}
On évaluera aussi, si le temps le permet, la généralisation du GNN à un taux
de bruit de test différent de celui de l'entraînement. Cette expérience
permet de mesurer une forme précise de robustesse au décalage de distribution ;
elle ne prouve pas à elle seule une robustesse à tous les modèles de bruit.

\paragraph{Portée honnête.}
Le projet est d'abord une reproduction contrôlée et une étude de sensibilité,
pas l'annonce d'un nouveau décodeur de l'état de l'art. Son intérêt dépendra
de la rigueur du protocole, de la qualité des mesures et de l'analyse des
conditions dans lesquelles le GNN se rapproche -- ou non -- de PyMatching.

% ============================================================
\section{0. Comprendre le problème physique et mathématique}
% ============================================================

\subsection{Pourquoi faut-il corriger des erreurs quantiques ?}

Un qubit est un système quantique dont l'état peut s'écrire
\[
  |\psi\rangle = \alpha |0\rangle+\beta |1\rangle,
  \qquad |\alpha|^2+|\beta|^2=1.
\]
Les amplitudes $\alpha$ et $\beta$ sont des nombres complexes. Une mesure
dans la base computationnelle donne $0$ ou $1$ avec des probabilités
$|\alpha|^2$ et $|\beta|^2$. Contrairement à un bit classique, l'information
quantique ne peut pas être copiée arbitrairement (théorème de non-clonage),
et une mesure directe du qubit logique risque de détruire l'information
qu'on souhaite préserver.

Les qubits physiques subissent des perturbations : erreurs de relaxation,
déphasage, erreurs de portes, erreurs de préparation et erreurs de mesure.
On utilise donc un \emph{code correcteur quantique} : l'information logique
est encodée collectivement dans plusieurs qubits physiques. Les mesures de
contrôle détectent des incompatibilités révélatrices d'erreurs sans mesurer
directement l'état logique.

\subsection{Qu'est-ce qu'une erreur physique ?}

Dans le formalisme de Pauli, les erreurs élémentaires usuelles sur un qubit
sont :
\[
 X=\begin{pmatrix}0&1\\1&0\end{pmatrix},\quad
 Z=\begin{pmatrix}1&0\\0&-1\end{pmatrix},\quad
 Y=iXZ=\begin{pmatrix}0&-i\\i&0\end{pmatrix}.
\]
\begin{itemize}
  \item $X$ échange $|0\rangle$ et $|1\rangle$ : on l'interprète comme une
  erreur de bit-flip.
  \item $Z$ change le signe relatif entre les amplitudes : c'est une erreur
  de phase-flip.
  \item $Y$ combine les deux effets, à une phase globale près.
\end{itemize}
Une erreur physique peut aussi être une erreur de mesure ou de porte dans
le circuit qui extrait le syndrome. Dans une simulation au niveau du circuit,
il ne faut donc pas imaginer que les seules erreurs se produisent sur les
qubits de données : des fautes peuvent survenir pendant les portes de
préparation, les portes d'extraction du syndrome, les opérations de remise
à zéro et les mesures.

\paragraph{Modèle dépolarisant.}
À titre d'exemple, un canal dépolarisant à un qubit peut être défini par
\[
 \mathcal D_p(\rho)=(1-p)\rho+
 \frac p3(X\rho X+Y\rho Y+Z\rho Z),
\]
où $\rho$ est la matrice densité de l'état et $p$ la probabilité qu'une
erreur de Pauli non triviale se produise. Les conventions exactes varient
selon l'emplacement de l'erreur et le simulateur. Le paramètre
\texttt{after\_clifford\_depolarization} de Stim ne signifie pas que toutes
les sources de bruit du circuit sont nécessairement réglées au même taux.
Le protocole doit donc déclarer explicitement chaque paramètre de bruit.

\subsection{Pourquoi ne pas mesurer directement si une erreur a eu lieu ?}

La mesure de l'état logique détruirait généralement l'information que la
mémoire doit préserver. Le code mesure à la place des observables de contrôle
appelées \emph{stabilisateurs}. Un stabilisateur est un opérateur de Pauli
$S_i$ tel que les états du sous-espace code satisfont
\[
 S_i|\psi_{\mathrm{code}}\rangle=+|\psi_{\mathrm{code}}\rangle.
\]
On mesure plusieurs stabilisateurs qui commutent entre eux. Chaque résultat
est $+1$ ou $-1$ (ou, de façon équivalente, un bit $0$ ou $1$ après
convention de codage). Ces mesures donnent des informations sur les erreurs
sans révéler directement l'état logique encodé.

\subsection{Qu'est-ce qu'un syndrome ? Qu'est-ce qu'un détecteur ?}

Il faut distinguer trois objets souvent confondus.

\paragraph{1. Résultat de mesure.}
Notons $m_{i,t}\in\{0,1\}$ le résultat binaire de la mesure du stabilisateur
numéro $i$ au cycle $t$, après conversion du résultat physique en bit.

\paragraph{2. Événement de détection.}
Un détecteur compare des résultats de mesure qui devraient être compatibles
en l'absence d'erreur. Dans le cas courant de mesures répétées d'un même
stabilisateur, on peut construire
\[
 d_{i,t}=m_{i,t}\xor m_{i,t-1}.
\]
Ici $\xor$ désigne le XOR, ou addition modulo 2 : $0\xor0=0$,
$0\xor1=1$, $1\xor0=1$ et $1\xor1=0$. Si le résultat ne change pas,
$d_{i,t}=0$ ; s'il change, $d_{i,t}=1$. Un détecteur vaut donc 1 lorsqu'une
relation attendue entre mesures est violée. Pour le premier cycle, la
relation peut comparer le résultat à une préparation connue ; au dernier
cycle, des relations particulières peuvent impliquer la mesure finale.
La formule ci-dessus est une intuition utile, mais la définition exacte des
détecteurs dépend du circuit et de ses annotations.

\paragraph{3. Syndrome.}
Le syndrome d'un tirage est le vecteur de tous les événements de détection :
\[
 \mathbf s=(d_1,d_2,\ldots,d_{N_D})\in\{0,1\}^{N_D},
\]
où $N_D$ est le nombre de détecteurs annotés dans le circuit. Le syndrome
n'est pas nécessairement un unique bit : il contient les résultats de tous
les détecteurs, souvent répartis dans l'espace et au cours du temps.

Un détecteur qui vaut 1 est souvent décrit comme \og activé \fg{} ou
\og allumé \fg{}. Cela ne veut pas dire qu'il identifie exactement le qubit
fautif. Plusieurs erreurs physiques différentes peuvent activer les mêmes
détecteurs ; certaines erreurs peuvent même ne pas produire de changement
de syndrome mais agir sur l'information logique. Le décodeur doit inférer
l'effet logique à partir de l'ensemble des détecteurs.

\subsection{Exemple simple : pourquoi les détecteurs sont-ils utiles ?}

Considérons un stabilisateur idéal mesuré trois fois. Sans erreur, supposons
que les résultats soient
\[
 (m_1,m_2,m_3)=(0,0,0).
\]
Les deux détecteurs temporels valent alors
\[
 (d_2,d_3)=(m_2\xor m_1,m_3\xor m_2)=(0,0).
\]
Si une erreur de mesure fait basculer uniquement la mesure centrale, on
observe par exemple $(0,1,0)$ et donc
\[
 (d_2,d_3)=(1,1).
\]
Deux détecteurs s'activent : l'incohérence apparaît à l'entrée et à la sortie
temporelles de l'erreur de mesure. En revanche, une erreur persistante sur
un qubit de données peut produire un motif différent, souvent spatialement
localisé. C'est pourquoi les détecteurs sont naturellement organisés dans
l'espace et le temps.

Cet exemple est schématique. Dans un vrai circuit quantique, les détecteurs
sont définis par les relations de parité exactes entre mesures, y compris
les effets des portes d'extraction du syndrome et des conditions aux limites.

\subsection{Pourquoi répéter les mesures dans le temps ?}

Une mesure de stabilisateur peut elle-même être fautive. Une seule série de
mesures ne permet donc pas toujours de distinguer une erreur de mesure d'une
erreur sur les qubits de données. On répète les mesures sur plusieurs cycles.
La structure spatio-temporelle du syndrome aide alors à différencier les
motifs d'erreur. Pour une mémoire de code de surface de distance $d$, on
choisit souvent un nombre de cycles proportionnel à $d$ pour un protocole
de mémoire standard, mais la valeur exacte dépend de l'expérience simulée.

\subsection{Qu'est-ce qu'une erreur logique ?}

Un code protège un qubit logique encodé dans plusieurs qubits physiques.
Le décodeur ne cherche pas nécessairement à identifier exactement chaque
erreur physique survenue. Il doit choisir une correction qui ramène l'état
dans le sous-espace code et préserve l'information logique.

Deux suites d'erreurs physiques peuvent produire le même syndrome, mais
avoir des effets logiques différents. Leur différence peut correspondre à
un opérateur logique non trivial. Une correction peut donc annuler les
défauts locaux tout en appliquant, sans le vouloir, une opération logique
$X_L$ ou $Z_L$ sur le qubit encodé. C'est ce qu'on appelle un échec logique.

Dans Stim, l'observable logique est suivi au moyen d'une annotation telle
que \texttt{OBSERVABLE\__INCLUDE}. Pour un tirage donné, le simulateur
fournit les événements de détection et le bit d'observable logique retourné
par l'expérience simulée. Ce bit constitue la cible de comparaison pour le
décodeur. Il faut respecter la convention exacte du circuit utilisé :
le bit observable de référence est le résultat logique effectivement observé
dans le tirage, et le décodeur prédit le retournement logique à corriger.

\subsection{Qu'est-ce que le code de surface et sa distance ?}

Le code de surface est un code correcteur quantique défini sur un réseau
bidimensionnel, avec des stabilisateurs locaux de type $X$ et $Z$. Les
interactions locales permettent une extraction du syndrome compatible avec
des architectures matérielles de qubits voisins.

La \emph{distance} $d$ du code est la longueur minimale, en poids de Pauli,
d'un opérateur logique non trivial. Pour un code de surface standard,
augmenter $d$ augmente la taille du code et le nombre d'erreurs physiques
nécessaires pour former une chaîne logique non triviale. Cela améliore la
protection lorsque le bruit se situe dans le régime approprié, mais augmente
aussi le coût matériel et computationnel. Les distances $d=3$ et $d=5$ sont
de petits cas adaptés à un prototype ; ils ne suffisent pas à établir des
conclusions d'échelle pour de grandes mémoires quantiques.

\subsection{Comment le décodeur agit-il en pratique ?}

Le circuit bruité est simulé ; chaque tirage (\emph{shot}) produit :
\begin{enumerate}
  \item un vecteur de détecteurs activés $\mathbf s$ ;
  \item une ou plusieurs valeurs de référence pour les observables logiques ;
  \item éventuellement, des informations supplémentaires pour analyser le
  tirage, mais qui ne doivent pas être données au décodeur si elles ne sont
  pas disponibles en expérience réelle.
\end{enumerate}

Le décodeur reçoit le syndrome $\mathbf s$ et prédit le retournement de
l'observable logique. On compare ensuite la prédiction à l'observable de
référence produite par le simulateur. Le décodeur n'a pas accès à la
description complète des fautes qui ont réellement eu lieu : ce serait une
fuite d'information et rendrait l'évaluation artificiellement facile.

\subsection{Les trois niveaux à ne pas confondre}

\begin{center}
\begin{tabularx}{\textwidth}{>{\bfseries}p{3.2cm} X}
\toprule
Niveau & Question \\
\midrule
Erreur physique & Quelle faute de porte, de préparation, de mesure ou de qubit s'est produite ? \\
Événements de détection & Quelles relations attendues entre les mesures sont violées ? \\
Erreur logique & Après la correction choisie, l'information logique est-elle incorrecte ? \\
\bottomrule
\end{tabularx}
\end{center}

Le décodeur résout le troisième problème en exploitant le deuxième ; il ne
reconstruit pas nécessairement le premier.

% ============================================================
\section{1. Problématique, hypothèses et périmètre}
% ============================================================

\subsection{Problématique principale}

Pour $d\in\{3,5\}$ et pour plusieurs taux physiques $p$, on compare le taux
d'erreur logique d'un GNN entraîné sur des exemples simulés à celui de
PyMatching, sur un ensemble de test commun et indépendant. On mesure aussi
le temps d'inférence et le temps d'entraînement.

\[
\boxed{\text{Quelle différence de qualité et de coût observe-t-on entre
le GNN et PyMatching dans un protocole contrôlé ?}}
\]

La question n'est pas \og le GNN est-il meilleur en général ?\fg{} mais
\og dans quelles conditions précises, sur quelle tâche et avec quelle
incertitude, le GNN atteint-il tel niveau de performance ?\fg{}

\subsection{Hypothèses expérimentales}

\begin{description}
 \item[H1 -- apprentissage.] Le GNN peut apprendre une relation utile entre
 le syndrome et l'observable logique à partir d'exemples simulés.
 \item[H2 -- performance.] Le taux d'erreur logique du GNN peut être proche
 de celui de PyMatching, mais ce résultat n'est pas garanti et doit être
 mesuré.
 \item[H3 -- coût.] L'entraînement du GNN nécessite un coût initial, tandis
 que le temps d'inférence peut être compétitif. Cette hypothèse dépend de
 l'architecture, du matériel et de la taille des lots.
 \item[H4 -- généralisation.] Un GNN entraîné à un taux de bruit donné peut
 voir ses performances varier lorsque le taux de bruit change. On mesurera
 cette variation sans supposer qu'elle sera nécessairement défavorable.
\end{description}

\subsection{Périmètre recommandé}

\begin{itemize}
 \item Code de surface en mémoire, avec une tâche Stim clairement identifiée.
 \item Distances $d=3$ et $d=5$.
 \item Bruit dépolarisant au niveau du circuit, avec tous les paramètres
 déclarés explicitement.
 \item PyMatching comme référence principale.
 \item Un GNN de petite taille, puis éventuellement une ablation simple.
 \item Données simulées pour l'expérience principale ; données Google
 uniquement en extension après validation du pipeline.
\end{itemize}

Il faut fixer la tâche exacte avant de commencer : par exemple
\texttt{surface\_code:rotated\_memory\_z} ou
\texttt{surface\_code:rotated\_memory\_x}. Les deux tâches ne doivent pas
être mélangées sans explication, car elles protègent des observables logiques
différentes.

% ============================================================
\section{2. Méthode et protocole expérimental}
% ============================================================

\subsection{Étape A : générer un circuit bruité avec Stim}

Stim est un simulateur optimisé pour les circuits de correction d'erreurs,
notamment les circuits de Clifford. Un exemple minimal est :

\begin{lstlisting}[language=Python,caption={Exemple de circuit de mémoire de code de surface}]
import stim

d = 3
p = 0.005

circuit = stim.Circuit.generated(
    "surface_code:rotated_memory_z",
    distance=d,
    rounds=d,
    after_clifford_depolarization=p,
    before_round_data_depolarization=p,
    before_measure_flip_probability=p,
    after_reset_flip_probability=p,
)
\end{lstlisting}

Ce code est un point de départ et non une définition universelle du
\og bruit dépolarisant \fg{}. Il règle ici plusieurs sources de bruit à la
même valeur pour créer un scénario explicite. Pour l'expérience finale,
documenter chaque paramètre et vérifier les hypothèses du modèle.

\subsection{Étape B : produire les exemples étiquetés}

On échantillonne les détecteurs et les observables logiques séparément :
\begin{lstlisting}[language=Python,caption={Échantillonnage des syndromes et observables}]
import numpy as np

sampler = circuit.compile_detector_sampler()
syndrome, observables = sampler.sample(
    shots=10000,
    separate_observables=True,
)

print(syndrome.shape)     # (nombre de tirs, nombre de détecteurs)
print(observables.shape)  # (nombre de tirs, nombre d'observables)
\end{lstlisting}

La matrice \texttt{syndrome} contient les entrées accessibles au décodeur.
La matrice \texttt{observables} contient les cibles logiques de référence.
Elles doivent rester séparées : donner les observables au GNN en entrée
constituerait une fuite de cible.

\subsection{Étape C : construire la référence PyMatching}

\begin{lstlisting}[language=Python,caption={Décodeur MWPM de référence}]
import pymatching

dem = circuit.detector_error_model(decompose_errors=True)
matching = pymatching.Matching.from_detector_error_model(dem)

prediction_pm = matching.decode_batch(syndrome)
\end{lstlisting}

Le modèle d'erreur de détecteurs (\emph{detector error model}, DEM) décrit
les mécanismes de fautes sous la forme de relations entre détecteurs et
observables. L'option \texttt{decompose\_errors=True} demande une
décomposition des mécanismes pour obtenir une représentation compatible
avec le graphe de PyMatching dans les cas pris en charge. Il faut vérifier
les avertissements de Stim : une décomposition n'est pas une garantie que
toute corrélation complexe est parfaitement représentée par un graphe
simple.

\subsection{Étape D : calculer le taux d'erreur logique}

Pour une observable logique, la métrique est :
\[
 \widehat{\pL}=\frac{1}{N}\sum_{i=1}^{N}
 \mathbf{1}\!\left[\hat o_i\ne o_i\right],
\]
où $N$ est le nombre de tirages, $o_i\in\{0,1\}$ l'observable de référence,
$\hat o_i$ la prédiction du décodeur et $\mathbf 1[\cdot]$ vaut 1 si la
condition est vraie et 0 sinon.

\begin{lstlisting}[language=Python,caption={Calcul de l'erreur logique}]
num_failures = np.count_nonzero(prediction_pm != observables)
logical_error_rate = num_failures / observables.size
print("Taux d'erreur logique :", logical_error_rate)
\end{lstlisting}

Pour une seule observable, cette formule est simple. Si plusieurs observables
sont prédites, il faut définir la métrique avant l'expérience. Pour le taux
de tirages contenant au moins une erreur logique :
\[
 \widehat{\pL}_{\mathrm{any}}=
 \frac1N\sum_{i=1}^N
 \mathbf{1}\!\left[\exists j:\hat o_{ij}\ne o_{ij}\right].
\]
Ne pas confondre cette quantité avec le taux d'erreurs calculé par bit
d'observable. Pour le premier projet, une tâche à une observable logique
évite cette ambiguïté.

\subsection{Étape E : construire les graphes pour le GNN}

Un graphe est un objet mathématique $G=(V,E)$ : $V$ est l'ensemble des
nœuds et $E$ l'ensemble des arêtes reliant certains nœuds. Dans un GNN,
chaque nœud porte des attributs numériques et les couches du réseau
échangent des informations le long des arêtes.

\paragraph{Choix de représentation.}
Une première idée consiste à créer un nœud uniquement pour chaque
détecteur activé. Cette représentation est compacte, mais elle peut donner
un graphe vide et perdre la distinction entre absence d'un événement et
absence d'un nœud. Une solution plus stable pour un premier prototype est
de conserver tous les détecteurs potentiels comme nœuds et de leur attribuer
un bit d'activation $d_i\in\{0,1\}$. On peut ajouter les coordonnées
spatio-temporelles $(x_i,y_i,t_i)$ lorsque le circuit les fournit.

Les arêtes ne doivent pas être choisies uniquement parce que deux points
sont proches en distance euclidienne. Elles devraient refléter une
structure pertinente du code ou des relations d'erreur. Le modèle de
détecteurs de Stim peut servir de référence pour comprendre quelles
combinaisons de détecteurs sont reliées par des mécanismes de faute. Une
construction par $k$ plus proches voisins reste possible comme variante,
mais elle doit être motivée et évaluée comme un choix de modèle.

\paragraph{Attention à la cible.}
Le GNN doit prédire l'observable logique de référence à partir du syndrome,
pas la localisation complète de l'erreur physique. En général, plusieurs
histoires d'erreurs peuvent être compatibles avec un même syndrome. La
cible est donc une classe logique, et non une vérité unique sur les fautes
microscopiques.

\subsection{Étape F : architecture initiale du GNN}

Une couche de passage de messages peut s'écrire de façon générique :
\[
 \mathbf h_i^{(\ell+1)}
 =\phi\!\left(
 \mathbf h_i^{(\ell)},
 \operatorname{AGG}_{j\in\mathcal N(i)}
 \psi(\mathbf h_i^{(\ell)},\mathbf h_j^{(\ell)},\mathbf e_{ij})
 \right),
\]
où :
\begin{itemize}
 \item $\mathbf h_i^{(\ell)}$ est le vecteur de caractéristiques du nœud $i$
 à la couche $\ell$ ;
 \item $\mathcal N(i)$ est l'ensemble de ses voisins ;
 \item $\mathbf e_{ij}$ représente éventuellement les attributs de l'arête ;
 \item $\psi$ calcule un message à partir du nœud et de son voisin ;
 \item $\operatorname{AGG}$ agrège les messages (somme, moyenne ou maximum) ;
 \item $\phi$ met à jour la représentation du nœud.
\end{itemize}
Après plusieurs couches, une opération de lecture globale rassemble les
représentations en un vecteur de graphe, par exemple
\[
 \mathbf h_G=\frac{1}{|V|}\sum_{i\in V}\mathbf h_i^{(L)}.
\]
Une tête de classification prédit ensuite une probabilité de retournement
logique :
\[
 \hat q=\sigma(f_\theta(\mathbf h_G)),\qquad
 \sigma(z)=\frac{1}{1+e^{-z}}.
\]
Le paramètre $\theta$ représente les poids appris du réseau. Le modèle peut
être entraîné avec la perte d'entropie croisée binaire :
\[
 \mathcal L(\theta)=-\frac1N\sum_{i=1}^{N}
 \left[y_i\log\hat q_i+(1-y_i)\log(1-\hat q_i)\right],
\]
où $y_i$ est l'étiquette binaire et $\hat q_i$ la probabilité prédite.
En pratique, \texttt{BCEWithLogitsLoss} est généralement préférable à
l'application manuelle du logarithme et de la sigmoïde, pour des raisons
numériques.

Ces équations décrivent une famille d'architectures. Pour prétendre
reproduire Lange et al., il faudra reprendre leur définition effective du
graphe, des caractéristiques, de la cible et de la procédure d'entraînement
plutôt que supposer que ce GNN générique est identique au leur.

\subsection{Étape G : séparation des données}

Séparer les tirages avant l'entraînement :
\begin{itemize}
 \item entraînement : ajuste les paramètres du GNN ;
 \item validation : choisit les hyperparamètres et l'arrêt de l'entraînement ;
 \item test : fournit le résultat final, consulté après les décisions de modèle.
\end{itemize}
Les tirages étant indépendants lorsqu'ils sont échantillonnés séparément du
même circuit, le risque principal de fuite est surtout une mauvaise
organisation du code ou l'utilisation accidentelle des étiquettes de test
pour choisir le modèle. Les mêmes échantillons de test doivent servir à
PyMatching et au GNN, afin que la comparaison soit appariée.

\subsection{Étape H : généralisation au changement de bruit}

On entraîne par exemple à $p_{\mathrm{train}}=0.005$, puis on teste à
$p_{\mathrm{test}}\in\{0.0025,0.005,0.01\}$, sous réserve que ces valeurs
soient adaptées au protocole choisi. Pour chaque valeur de test, on génère
un circuit correspondant. PyMatching doit être construit à partir du modèle
de bruit de ce circuit de test, tandis que le GNN est conservé sans
réentraînement pour mesurer son transfert.

Il s'agit d'une comparaison volontairement asymétrique :
PyMatching reçoit le modèle de bruit correct au test, le GNN non. C'est un
test de robustesse d'un décodeur appris, pas une comparaison où les deux
méthodes disposent des mêmes informations a priori.

% ============================================================
\section{3. Métriques et analyse statistique}
% ============================================================

\subsection{Métrique principale : taux d'erreur logique}

La métrique centrale est le taux d'erreur logique $\pL$. Il mesure la
fraction des tirages pour lesquels le décodeur ne prédit pas correctement
l'observable logique. Une accuracy élevée peut être trompeuse lorsque les
échecs sont rares ; il faut donc rapporter directement le taux d'erreur
logique et le nombre absolu d'échecs.

\subsection{Différence entre les décodeurs}

Pour un même point expérimental, on rapporte la différence absolue
\[
 \Delta \pL=\pL^{\GNN}-\pL^{\PM}
\]
et, lorsque $\pL^{\PM}$ n'est pas trop proche de zéro, l'écart relatif
\[
 \delta_{\mathrm{rel}}=
 \frac{\pL^{\GNN}-\pL^{\PM}}{\pL^{\PM}}.
\]
Un écart relatif positif signifie que le GNN échoue plus souvent ; un écart
négatif signifie qu'il échoue moins souvent. Si le dénominateur est nul ou
très faible, l'écart relatif est instable : présenter alors la différence
absolue et les intervalles de confiance.

\subsection{Incertitude statistique}

Si $K$ échecs sont observés sur $N$ tirages indépendants, l'estimateur est
$\hat p=K/N$. Pour une première analyse, un intervalle de confiance binomial
à 95\,\% (par exemple Wilson ou Clopper--Pearson) est préférable à une
approximation gaussienne lorsque $K$ est petit.

Si aucun échec n'est observé, cela ne prouve pas que le vrai taux est nul.
Une règle approximative courante donne une borne supérieure à 95\,\%
d'environ $3/N$ dans ce cas (\og règle de trois \fg{}). Ainsi, zéro échec
sur 1\,000 tirs indique seulement que le taux est compatible avec une valeur
de l'ordre de quelques $10^{-3}$ ou moins ; il ne permet pas de mesurer
fiablement un taux de $10^{-5}$.

Un objectif pratique est de viser environ 100 échecs observés par point de
mesure lorsque le coût de simulation le permet. Le nombre de tirs nécessaire
dépend du taux réel : si le taux est très faible, il faudra beaucoup plus
de tirs. Il est donc utile de commencer par une courte simulation pilote.

\subsection{Temps de calcul}

Mesurer séparément :
\begin{enumerate}
 \item temps de génération des échantillons ;
 \item temps de construction du DEM et de PyMatching ;
 \item temps d'entraînement du GNN ;
 \item temps d'inférence, après échauffement éventuel, sur un nombre de
 tirages annoncé.
\end{enumerate}
Ne pas mélanger le temps de génération du circuit au temps de décodage.
Préciser le processeur, le GPU éventuel, les tailles de lots et le nombre
de répétitions du benchmark. Le débit peut être exprimé en tirs par seconde.
Une comparaison du temps d'inférence doit être interprétée avec prudence :
PyMatching et un GNN peuvent ne pas utiliser le même matériel ni la même
stratégie de traitement par lots.

\subsection{Tableau de résultats à remplir}

\begin{center}
\begin{tabularx}{\textwidth}{c c c c c c}
\toprule
$d$ & $p$ & $N$ & Échecs GNN & Échecs PyMatching & Temps GNN / PM \\
\midrule
3 & & & & & \\
3 & & & & & \\
5 & & & & & \\
5 & & & & & \\
\bottomrule
\end{tabularx}
\end{center}

Pour chaque ligne, ajouter dans le rapport les taux calculés, leurs
intervalles de confiance, les graines utilisées et la configuration exacte
du bruit.

% ============================================================
\section{4. État de l'art et références commentées}
% ============================================================

\subsection{Lange et al. : le travail directement étudié}

\paragraph{Référence.}
M. Lange et al., \emph{Data-driven decoding of quantum error correcting
codes using graph neural networks}, \emph{Physical Review Research},
\textbf{7}, 023181 (2025), version arXiv:2307.01241, version révisée en
mars 2025. DOI : \url{https://doi.org/10.1103/PhysRevResearch.7.023181}.

\paragraph{Problème posé.}
Les décodeurs optimaux au sens de la vraisemblance peuvent être coûteux.
Les décodeurs efficaces, tels que le matching, utilisent une structure
algorithmique et un modèle d'erreur. Leur qualité peut dépendre de la
justesse de ce modèle et de la capacité de la représentation à capturer les
corrélations. Lange et ses coauteurs étudient une autre voie : apprendre
directement, à partir d'exemples, à associer une configuration de détecteurs
à une classe d'erreur logique.

\paragraph{Formulation.}
Le syndrome est transformé en un graphe annoté. Le GNN agrège l'information
locale et produit une prédiction de la classe logique. Le modèle apprend à
partir de données simulées ; les observables logiques correspondantes
servent de cibles d'entraînement. L'approche est dite \emph{data-driven}
car elle apprend une relation entrée-sortie au lieu de calculer uniquement
une solution par optimisation sur le graphe de matching.

\paragraph{Résultats et interprétation.}
Les auteurs rapportent des performances compétitives, et dans certains
régimes supérieures au matching pour le bruit au niveau du circuit sur des
données simulées, alors même que le décodeur de matching dispose du modèle
de bruit. Ils étudient également des données expérimentales de Google pour
le code de répétition. La version révisée étend les résultats simulés à des
distances allant jusqu'à $d=9$ et ajoute une comparaison avec
\emph{belief-matching}. Ces résultats sont ceux de leur protocole, de leur
architecture et de leurs données ; ils ne garantissent pas que tout GNN
simple obtiendra la même performance.

\paragraph{Ce qu'il faut lire attentivement.}
\begin{itemize}
 \item La représentation du graphe et les caractéristiques des nœuds.
 \item La définition de la sortie du réseau et des classes logiques.
 \item Les détails du modèle de bruit et la façon de générer les données.
 \item La séparation entraînement/validation/test et le nombre d'échantillons.
 \item Les courbes de taux d'erreur logique, pas seulement l'accuracy.
 \item Le coût d'entraînement, le temps d'inférence et les limites d'échelle.
\end{itemize}

\paragraph{Rôle dans notre projet.}
C'est la référence méthodologique principale. Pour parler de reproduction,
il faudra reprendre une configuration identifiable de l'article et comparer
les résultats aux figures et paramètres correspondants. Si l'on change
simultanément le graphe, l'architecture, la cible et le bruit, le projet
devient une implémentation inspirée de l'article plutôt qu'une reproduction
fidèle.

\subsection{Higgott : PyMatching et l'appariement de poids minimal}

\paragraph{Référence.}
O. Higgott, \emph{PyMatching: A Python package for decoding quantum codes
with minimum-weight perfect matching}, arXiv:2105.13082 (2021).
Documentation : \url{https://pymatching.readthedocs.io/en/stable/}.

\paragraph{Idée mathématique du matching.}
Le décodeur construit un graphe dont les sommets correspondent aux
détecteurs. Les mécanismes d'erreur qui provoquent un ou deux événements
peuvent être représentés par des arêtes, éventuellement vers une frontière.
Chaque arête porte un poids lié à la probabilité de l'événement d'erreur.
Sous un modèle simple, un poids usuel est le logarithme des cotes :
\[
 w_e=\log\frac{1-q_e}{q_e},
\]
où $q_e$ est la probabilité associée à l'événement représenté par l'arête.
La convention précise des poids dépend de la construction du modèle. Le
décodeur cherche un ensemble de chemins ou d'arêtes qui explique le
syndrome avec un coût total minimal, sous les contraintes du problème
d'appariement.

Le poids logarithmique se comprend par la probabilité d'un ensemble
d'événements indépendants. Si $E$ est l'ensemble des événements choisis,
alors, à un facteur commun près,
\[
 -\log \Prb(E)
 =\sum_{e\in E}-\log q_e
  +\sum_{e\notin E}-\log(1-q_e).
\]
En retirant la somme constante des termes $-\log(1-q_e)$, on obtient un
coût proportionnel à
\[
 \sum_{e\in E}\log\frac{1-q_e}{q_e}.
\]
Minimiser ce coût revient, dans ce modèle, à rechercher l'explication la
plus probable parmi les configurations représentées. Cette interprétation
a des limites : les événements peuvent être corrélés, plusieurs configurations
peuvent produire le même syndrome, et le minimum de poids n'est pas toujours
la décision optimale au sens de la probabilité totale d'une classe logique.

\paragraph{Forces et limites.}
PyMatching est rapide et constitue une référence robuste. Il utilise
explicitement le modèle de bruit converti en DEM. Un modèle incorrect ou
une représentation qui ignore certaines corrélations peut dégrader ses
résultats. Inversement, il serait injuste de lui comparer un GNN en lui
fournissant volontairement un modèle de bruit faux : la référence principale
doit utiliser le modèle correspondant au circuit simulé.

\paragraph{Rôle dans notre projet.}
PyMatching fournit la référence quantitative principale. Le protocole doit
le faire décoder exactement les mêmes syndromes que le GNN, avec la même
définition d'échec logique.

\subsection{Google Quantum AI : le contexte expérimental}

\paragraph{Référence.}
Google Quantum AI, \emph{Suppressing quantum errors by scaling a surface
code logical qubit}, \emph{Nature} \textbf{614}, 676--681 (2023).
DOI : \url{https://doi.org/10.1038/s41586-022-05434-1}.
Données : \url{https://zenodo.org/records/6804040}.

\paragraph{Pourquoi ce travail est important.}
Cette étude expérimentale montre comment la performance d'une mémoire
quantique encodée dans un code de surface évolue lorsque la distance du code
augmente. Elle est une référence importante pour comprendre la correction
d'erreurs sur du matériel réel, où les erreurs ne suivent pas nécessairement
un modèle dépolarisant indépendant et homogène.

\paragraph{Ce que les données ne garantissent pas.}
Un jeu de données expérimental ne constitue pas automatiquement un ensemble
prêt à l'emploi pour entraîner n'importe quel GNN. Il faut vérifier le
format des échantillons, les conventions des observables, les décodeurs
déjà appliqués, la séparation entre données d'entraînement et de test, et
la correspondance entre la tâche publiée et la tâche simulée. Les données
peuvent permettre une extension intéressante, mais elles ne sont pas
nécessaires pour valider le premier pipeline.

\subsection{Mihailov et al. : une autre approche neuronale}

\paragraph{Référence.}
A. Mihailov et al., \emph{Neural network decoder for near-term surface-code
experiments}, \emph{Physical Review Research} (2025).
DOI : \url{https://doi.org/10.1103/PhysRevResearch.7.013029}.

\paragraph{Intérêt de cette référence.}
Elle élargit le paysage au-delà d'un seul choix de GNN et rappelle que
l'efficacité d'un décodeur neuronal dépend de son architecture, de la
structure des données et du régime expérimental. Il faut lire l'article
pour identifier précisément la tâche, les architectures et les conditions
de comparaison avant de lui attribuer une conclusion plus spécifique.

\paragraph{Rôle dans notre projet.}
Lecture complémentaire utile pour la discussion finale : elle aide à
expliquer pourquoi \og réseau de neurones \fg{} n'est pas une méthode unique,
et pourquoi une comparaison propre doit décrire le modèle exact.

\subsection{Stim : outil de simulation et convention des détecteurs}

\paragraph{Référence logicielle.}
C. Gidney, \emph{Stim: a fast stabilizer circuit simulator}, Quantum 5,
497 (2021), DOI : \url{https://doi.org/10.22331/q-2021-07-06-497}.
Code et documentation : \url{https://github.com/quantumlib/Stim}.

Stim génère des circuits bruités et échantillonne rapidement les événements
de détection et les observables logiques. Les instructions de détecteur
(\texttt{DETECTOR}) indiquent les relations de parité dont la violation
signale un événement. Le DEM condense les mécanismes de bruit en relations
entre événements de détection et observables. Il est crucial de comprendre
que les détecteurs sont des annotations mathématiques liées au circuit,
et non simplement des capteurs physiques qui identifient directement
l'emplacement de chaque erreur.

\subsection{Synthèse comparative}

\begin{center}
\small
\begin{tabularx}{\textwidth}{p{2.5cm} X X}
\toprule
Méthode / travail & Information exploitée & Rôle dans l'étude \\
\midrule
PyMatching / MWPM & Syndrome + graphe d'erreurs pondéré issu du modèle de bruit
& Référence algorithmique rapide \\
Lange et al. (GNN) & Exemples syndrome--observable et graphe annoté
& Référence méthodologique à reproduire \\
Google et al. & Mesures d'une mémoire quantique physique et analyses de décodeurs
& Contexte réel, extension optionnelle \\
Mihailov et al. & Modèle neuronal pour le décodage du code de surface
& Alternative pour élargir l'état de l'art \\
Stim & Circuit bruité, détecteurs, observables et DEM
& Génération des données et définition de la tâche \\
\bottomrule
\end{tabularx}
\end{center}

% ============================================================
\section{5. Que peut apporter ce projet au-delà d'une reproduction ?}
% ============================================================

\subsection{Réponse honnête}

Si le projet consiste uniquement à exécuter le code des auteurs sur une
configuration déjà publiée et à reproduire une courbe, son apport
scientifique nouveau est faible. C'est une reproduction utile pour apprendre
et vérifier la compréhension, mais il ne faut pas la présenter comme une
nouvelle méthode ou comme une avancée de l'état de l'art.

En revanche, une reproduction bien documentée peut devenir une étude
expérimentale intéressante si elle répond à une question supplémentaire
qui n'est pas déjà traitée de la même manière dans l'article. La valeur
ajoutée doit être délimitée et vérifiée après lecture attentive de l'article.

\subsection{Extension recommandée : généralisation au taux de bruit}

L'extension la plus réaliste est de séparer deux questions :
\begin{enumerate}
 \item \textbf{Régime interpolé :} le GNN est entraîné et testé sous le même
 modèle de bruit, avec des tirages indépendants.
 \item \textbf{Régime décalé :} le GNN est entraîné à $p_{\mathrm{train}}$,
 puis testé à des valeurs $p_{\mathrm{test}}\ne p_{\mathrm{train}}$ sans
 réentraînement.
\end{enumerate}

Pour chaque valeur de $p_{\mathrm{test}}$, PyMatching est configuré à partir
du modèle de bruit correspondant au test. On mesure
\[
 \Delta\pL(p_{\mathrm{test}})=
 \pL^{\GNN}_{\mathrm{test}}(p_{\mathrm{test}})
 -\pL^{\PM}_{\mathrm{test}}(p_{\mathrm{test}}).
\]
Cette courbe indique si l'écart entre les méthodes évolue lorsque le bruit
change. Il faut préciser qu'il s'agit d'un test de généralisation à un
paramètre, et non d'une étude complète de robustesse à tous les défauts
réalistes (bruit biaisé, dérive temporelle, erreurs corrélées, bruit de
lecture non uniforme, etc.).

\subsection{Extension alternative : ablation de la représentation du graphe}

Si le pipeline de base fonctionne, on peut comparer deux représentations :
\begin{itemize}
 \item \textbf{Graphe physique / DEM :} les connexions sont motivées par
 les relations entre détecteurs et mécanismes de faute du circuit.
 \item \textbf{Graphe de proximité :} les connexions sont déterminées par
 les $k$ plus proches voisins des coordonnées des détecteurs.
\end{itemize}
On garde autant que possible le même modèle, la même procédure
d'entraînement, les mêmes tirages et le même budget d'optimisation. On
mesure si la structure d'arêtes change le taux d'erreur logique ou le temps
d'inférence. Cette ablation est plus informative qu'un choix arbitraire de
$k$, mais elle exige du temps et une définition rigoureuse des deux graphes.

\subsection{Extension expérimentale Google : intérêt et limites}

Des données réelles permettraient d'examiner le décalage entre une simulation
idéalisée et un appareil quantique. Cependant, c'est une extension coûteuse :
il faut aligner les conventions de détecteurs et d'observables, identifier
les partitions de données et vérifier ce qui est réellement disponible dans
l'archive. Elle ne doit être entreprise qu'après validation de bout en bout
du pipeline simulé.

\subsection{Quel apport peut-on revendiquer ?}

\begin{center}
\begin{tabularx}{\textwidth}{p{3.4cm} X}
\toprule
Réalisation & Revendication raisonnable \\
\midrule
Exécution du code publié sans changement & Reproduction technique et apprentissage ; pas de nouveauté méthodologique. \\
Réimplémentation indépendante et validation sur $d=3,5$ & Vérification reproductible d'un protocole sur un périmètre explicite. \\
Comparaison contrôlée entraînement/test à taux identique et différent & Étude de sensibilité au décalage de taux de bruit, si elle n'est pas déjà couverte de façon équivalente par la référence choisie. \\
Ablation graphe DEM contre proximité & Analyse de l'influence de la représentation, si le protocole est suffisamment contrôlé. \\
Test sur données expérimentales correctement alignées & Évaluation complémentaire sur des données réelles ; conclusion limitée au jeu de données disponible. \\
\bottomrule
\end{tabularx}
\end{center}

La formulation finale de l'apport doit dépendre de ce que révèle la lecture
de l'article. Si l'étude de généralisation ou l'ablation est déjà faite par
Lange et al. dans des conditions comparables, il faudra soit reproduire
explicitement cette partie, soit trouver une différence précise et défendable
dans le protocole. Il ne faut pas revendiquer une nouveauté simplement
parce qu'une expérience est nouvelle pour l'étudiant.

% ============================================================
\section{6. Planning opérationnel de 20 heures}
% ============================================================

\begin{center}
\begin{tabularx}{\textwidth}{X c}
\toprule
\textbf{Tâche} & \textbf{Durée} \\
\midrule
Lecture ciblée de Lange, Stim et PyMatching ; définition du protocole & 1 h \\
Pipeline Stim $\rightarrow$ PyMatching et validation du taux logique & 2 h \\
Référence statistique, décodeur trivial et préparation du test & 2 h \\
Construction et vérification des graphes & 3 h \\
Implémentation, entraînement et validation du GNN & 5 h \\
Comparaison des taux logiques et mesures de temps & 3 h \\
Test à un taux de bruit différent, si le pipeline est stable & 2 h \\
Rédaction, figures et limites & 2 h \\
\midrule
\textbf{Total} & \textbf{20 h} \\
\bottomrule
\end{tabularx}
\end{center}

\subsection{Critères de réussite}

Le projet est réussi si :
\begin{enumerate}
 \item le circuit et les paramètres de bruit sont documentés ;
 \item le calcul de l'échec logique est validé sur un exemple contrôlé ;
 \item PyMatching est évalué sur un jeu de test clairement défini ;
 \item le GNN apprend sur les données d'entraînement et est évalué sur un
 test indépendant ;
 \item les taux d'erreur sont accompagnés du nombre de tirs et d'une
 incertitude statistique ;
 \item le résultat négatif éventuel est analysé plutôt que masqué ;
 \item le code, les graines, les versions de bibliothèques et les paramètres
 sont consignés.
\end{enumerate}

\subsection{Risques et réponses prévues}

\begin{longtable}{p{4cm} p{5.5cm} p{5cm}}
\toprule
Risque & Conséquence & Réponse \\
\midrule
Taux d'échec très faible & Estimation statistique imprécise & Augmenter les tirs, commencer par un pilote et rapporter une borne plutôt que conclure à zéro. \\
GNN n'apprend pas & Pas de comparaison significative & Vérifier la cible, le déséquilibre des classes, la représentation du graphe et la fuite de données. \\
Graphe vide ou mal construit & Exemples perdus ou représentations incohérentes & Comparer graphe de détecteurs activés et graphe avec tous les détecteurs potentiels. \\
Temps d'entraînement trop long & Dépassement du budget & Réduire la taille du jeu, la dimension cachée ou le nombre de couches pour le prototype ; ne pas changer le test final. \\
Résultat différent de l'article & Impossible de conclure directement à un échec de reproduction & Vérifier le circuit, les paramètres de bruit, les cibles, les splits, l'architecture et les hyperparamètres. \\
Données expérimentales difficiles à aligner & Extension trop coûteuse & Garder l'extension Google hors du périmètre principal. \\
\bottomrule
\end{longtable}

% ============================================================
\section{7. Reproductibilité et bonnes pratiques}
% ============================================================

Le dépôt du projet devrait contenir :
\begin{itemize}
 \item un fichier de dépendances avec les versions de Python, Stim,
 PyMatching, PyTorch et PyTorch Geometric ;
 \item un script de génération des circuits et des données ;
 \item un script de construction du modèle PyMatching ;
 \item un script d'entraînement du GNN ;
 \item un script d'évaluation commun aux deux décodeurs ;
 \item un fichier de configuration listant $d$, les rounds, tous les taux
 de bruit, les graines, les tailles des ensembles et les hyperparamètres ;
 \item les résultats bruts (nombre de tirs, nombre d'échecs, durées) et les
 scripts qui produisent les figures.
\end{itemize}

Ne pas conserver uniquement les courbes finales : enregistrer également
les comptes d'échecs et le nombre de tirs permet de recalculer les taux et
les intervalles de confiance. Les résultats doivent pouvoir être régénérés
à partir des paramètres déclarés.

% ============================================================
\section{8. Références bibliographiques}
% ============================================================

\begin{thebibliography}{99}

\bibitem{lange2025}
M. Lange, P. Havström, B. Srivastava, I. Bengtsson, V. Bergentall,
K. Hammar, O. Heuts, E. van Nieuwenburg et M. Granath,
\emph{Data-driven decoding of quantum error correcting codes using graph
neural networks}, Physical Review Research \textbf{7}, 023181 (2025).
DOI : \url{https://doi.org/10.1103/PhysRevResearch.7.023181}.
Prépublication : \url{https://arxiv.org/abs/2307.01241}.

\bibitem{higgott2021}
O. Higgott, \emph{PyMatching: A Python package for decoding quantum codes
with minimum-weight perfect matching}, arXiv:2105.13082 (2021).
\url{https://arxiv.org/abs/2105.13082}.

\bibitem{gidney2021}
C. Gidney, \emph{Stim: a fast stabilizer circuit simulator},
Quantum \textbf{5}, 497 (2021).
DOI : \url{https://doi.org/10.22331/q-2021-07-06-497}.
Code : \url{https://github.com/quantumlib/Stim}.

\bibitem{google2023}
Google Quantum AI, \emph{Suppressing quantum errors by scaling a surface
code logical qubit}, Nature \textbf{614}, 676--681 (2023).
DOI : \url{https://doi.org/10.1038/s41586-022-05434-1}.
Données : \url{https://zenodo.org/records/6804040}.

\bibitem{mihailov2025}
A. Mihailov et al., \emph{Neural network decoder for near-term surface-code
experiments}, Physical Review Research (2025).
DOI : \url{https://doi.org/10.1103/PhysRevResearch.7.013029}.

\bibitem{stimdocs}
Quantumlib, \emph{Stim documentation: detector error model and sampling}.
\url{https://github.com/quantumlib/Stim}.

\bibitem{pymatchingdocs}
O. Higgott et al., \emph{PyMatching documentation}.
\url{https://pymatching.readthedocs.io/en/stable/}.

\bibitem{pygdocs}
PyTorch Geometric contributors, \emph{PyTorch Geometric documentation}.
\url{https://pytorch-geometric.readthedocs.io/}.

\end{thebibliography}

\end{document}
